\section{Background}
\label{sec:background}

\paragraph{Chunks and embeddings.}
Documents are cut into \emph{chunks}; here a chunk is an abstract's title
plus whole sentences, at most 512 tokens. An embedding model maps each
chunk and question to a vector so that similar texts are close. Dense
retrievers rank passages by the inner product of these
vectors~\cite{karpukhin2020dpr}. We use
\texttt{bge-small-en-v1.5}~\cite{xiao2024cpack} (384 dimensions) with unit
vectors, so the similarity is $s(q,c) = \langle e(q), e(c) \rangle$.

\paragraph{Top-$k$ retrieval.}
The retriever returns the $k$ chunks of corpus $C$ with the highest $s(q,c)$,
written $\mathrm{top}_k(q,C)$. A question with correct chunks $R(q)$
\emph{fails} when $R(q) \cap \mathrm{top}_k(q,C) = \emptyset$. We use $k=5$.

\paragraph{How an ingestion breaks a question.}
An ingestion adds chunks $N$. Let $m(q) = \max_{r \in R(q)} s(q,r)$. A
question that passed fails afterwards exactly when
\[
  \bigl|\{ n \in N : s(q,n) > m(q) \}\bigr|
  \;\ge\; k - \bigl|\{ c \in C \setminus R(q) : s(q,c) > m(q) \}\bigr| ,
\]
with ties going to older chunks. The correct chunk is not deleted but
\emph{displaced} by new chunks closer to the question. Since
$\|e(q)-e(n)\|^2 = 2 - 2\,s(q,n)$ for unit vectors, only questions near the
new chunks in the embedding space can break.

\paragraph{Distractors.}
Chunks that look relevant but do not answer the question are
\emph{distractors}. A few optimised passages are retrieved for many unseen
questions~\cite{zhong2023poisoning}, and five injected texts per question
can control the answer of a RAG system~\cite{zou2025poisonedrag}. We study
only displacement, which can be measured without a language model.
